% Figure F6 --- the harm-forensics loop (Section on forensics).
% \resizebox to \textwidth in main.tex.
\begin{tikzpicture}[
  st/.style={draw, rounded corners, align=center, minimum height=1.15cm,
             text width=2.05cm, font=\scriptsize, fill=black!3},
  arr/.style={-{Latex[length=2mm]}, thick, black!70},
  fb/.style={-{Latex[length=1.8mm]}, dashed, black!55},
  lbl/.style={font=\scriptsize\itshape, text=black!55},
  node distance=6mm
]
\node[st] (s1) {\textbf{Harm incident}\\an AI output or action causes harm};
\node[st, right=of s1] (s2) {\textbf{Reconstruct}\\trigger, conditions, taxonomy tag};
\node[st, right=of s2] (s3) {\textbf{Attribute}\\what the human did vs what the AI did};
\node[st, right=of s3] (s4) {\textbf{Scope \& cascade}\\downstream reach; who relied};
\node[st, right=of s4] (s5) {\textbf{Guardrail}\\matched to cause; dated boundary};
\node[st, right=of s5] (s6) {\textbf{Prevention}\\efficacy measured; recurrence reduced};
\draw[arr] (s1) -- (s2);
\draw[arr] (s2) -- (s3);
\draw[arr] (s3) -- (s4);
\draw[arr] (s4) -- (s5);
\draw[arr] (s5) -- (s6);
\draw[fb] (s6.south) to[out=-90,in=-90,looseness=0.3]
     node[below,lbl]{lowers the incidence of the failure class} (s1.south);
\node[lbl] at ($(s1.north)!0.5!(s6.north)+(0,6mm)$)
     {each step is a query over the immutable, linked record (Section~\ref{sec:datamodel})};
\end{tikzpicture}
